% Options for packages loaded elsewhere
\PassOptionsToPackage{unicode}{hyperref}
\PassOptionsToPackage{hyphens}{url}
%
\documentclass[
  10pt,
]{article}
\usepackage{amsmath,amssymb}
\usepackage{iftex}
\ifPDFTeX
  \usepackage[T1]{fontenc}
  \usepackage[utf8]{inputenc}
  \usepackage{textcomp} % provide euro and other symbols
\else % if luatex or xetex
  \usepackage{unicode-math} % this also loads fontspec
  \defaultfontfeatures{Scale=MatchLowercase}
  \defaultfontfeatures[\rmfamily]{Ligatures=TeX,Scale=1}
\fi
\usepackage{lmodern}
\ifPDFTeX\else
  % xetex/luatex font selection
\fi
% Use upquote if available, for straight quotes in verbatim environments
\IfFileExists{upquote.sty}{\usepackage{upquote}}{}
\IfFileExists{microtype.sty}{% use microtype if available
  \usepackage[]{microtype}
  \UseMicrotypeSet[protrusion]{basicmath} % disable protrusion for tt fonts
}{}
\makeatletter
\@ifundefined{KOMAClassName}{% if non-KOMA class
  \IfFileExists{parskip.sty}{%
    \usepackage{parskip}
  }{% else
    \setlength{\parindent}{0pt}
    \setlength{\parskip}{6pt plus 2pt minus 1pt}}
}{% if KOMA class
  \KOMAoptions{parskip=half}}
\makeatother
\usepackage{xcolor}
\usepackage[margin=0.85in]{geometry}
\setlength{\emergencystretch}{3em} % prevent overfull lines
\providecommand{\tightlist}{%
  \setlength{\itemsep}{0pt}\setlength{\parskip}{0pt}}
\setcounter{secnumdepth}{-\maxdimen} % remove section numbering
\usepackage{amsmath,amssymb,amsthm}
\DeclareUnicodeCharacter{03A3}{\ensuremath{\Sigma}}
\DeclareUnicodeCharacter{2212}{\ensuremath{-}}
\DeclareUnicodeCharacter{221E}{\ensuremath{\infty}}
\DeclareUnicodeCharacter{222A}{\ensuremath{\cup}}
\DeclareUnicodeCharacter{2265}{\ensuremath{\ge}}
\DeclareUnicodeCharacter{FB00}{ff}
\DeclareUnicodeCharacter{FB01}{fi}
\DeclareUnicodeCharacter{FB03}{ffi}
\ifLuaTeX
  \usepackage{selnolig}  % disable illegal ligatures
\fi
\usepackage{bookmark}
\IfFileExists{xurl.sty}{\usepackage{xurl}}{} % add URL line breaks if available
\urlstyle{same}
\hypersetup{
  pdftitle={Original-witness trace and a conditional countable Runge chain},
  pdfauthor={Unverified mathematical proof draft},
  hidelinks,
  pdfcreator={LaTeX via pandoc}}

\title{Original-witness trace and a conditional countable Runge chain}
\author{Unverified mathematical proof draft}
\date{5 October 2026}

\begin{document}
\maketitle

Candidate=false. The exact unchanged original P16.1 is not solved. The
results below retain every contact and geometric antecedent of the
task-granted round77 theorem and precise round80 PASS. The mathematical
deductions are drafts; documentary execution does not validate them.

\section{n1. Whole-base trace on the original
witness}\label{n1.-whole-base-trace-on-the-original-witness}

This is a conditional mathematical proof draft. The field \(k\) is
algebraically closed, trivially valued, and of characteristic zero,
without a cardinality restriction. Let \(S=D(a)\subset\mathbb A_k^n\),
let \(Y=\operatorname{Spec} B\) be integral, and let \(\pi:Y\to S\) be
finite and surjective, of function-field degree \(d\). Let
\(u:Y^{\rm an}\to\mathbb R\cup\{-\infty\}\) be usc, locally bounded
above, generically finite, and satisfy \(u\le Cg\circ\pi+B_0\) for an
actual affine FS exhaustion \(g\) of the original base ring. There is
ONE proper birational witness \(\mu:P\to Y\); on it the pullback is
locally a decreasing limit on a common neighborhood of the exact
rational contact weights in the task-granted round77/80 theorem.
Supports are in logarithms of units, lie below their weight, and agree
at the tested point with the same additive constant.

Retain the geometric antecedents of that accepted smooth gate in EVERY
smooth-open/base-removal application below: a smooth ample-boundary
projective compactification and projective birational domination of the
restricted original witness by a smooth projective source, with the
specified matching reference. This note does not prove those
compactification/domination/resolution antecedents. Write GEO for this
collection of applications, not for an arbitrary completed-residue-field
regularization theorem.

The task grants the entire positive norm-law theorem, not merely a
flat-fiber formula. With \(A=k[S]\) normal, its generic degree-\(d\)
norm and characteristic coefficients extend as polynomial laws over
\(A\). After specialization and algebraic closure, the norm is the
product of residue coordinates with strictly positive integer exponents.
Nilpotent units have trivial character; Cayley--Hamilton excludes a zero
exponent. For an original analytic fiber point \(y\), its weight is \[
 m_y=e_y[\kappa_y:\mathcal H(s)]>0,\qquad\sum_{y\mid s}m_y=d.
\] Here \(e_y\) is the residue-character exponent and \(\kappa_y\) the
finite residue field of the corresponding Artinian factor. These are not
nonflat fiber lengths. The weights regroup under every complete
extension. This granted result includes continuity of traces of
continuous functions on controlled compact preimages and local constancy
of the integer trace of a clopen indicator.

For clarity the granted finite test identity is \[
 \sum_{y\mid s}m_y D^{-1}\max_i(\log|b_i(y)|+\lambda_i)
 =D^{-1}\max_{|\alpha|=d}\bigl(\log|a_\alpha(s)|+\alpha\cdot\lambda\bigr),
 \quad N\Bigl(\sum_iT_i b_i\Bigr)=\sum_\alpha a_\alpha T^\alpha.
 \tag{T1}
\] The constant branch is retained, and the original base generators can
be appended at dominated positive slopes. Weighted Gauss
multiplicativity applies at every positive real radius. The theorem
proved here under GEO is that the left trace of the ORIGINAL weak \(u\),
with every finite or minus-infinity value, is an actual decreasing
original-ring affine FS limit on the WHOLE base.

\subsection{Construct the base extension before identifying the bad
values}\label{construct-the-base-extension-before-identifying-the-bad-values}

Choose one principal \(W=D(f)\subset S\) on which \(\pi\) is finite
étale and the original witness is an isomorphism over \(\pi^{-1}W\).
Such an open exists: in characteristic zero the generic extension is
separable; clear denominators in a primitive-element presentation, then
invert its discriminant and the finitely many remaining module
denominators. Remove also the finite image of the complement of the
witness's isomorphism locus. A nonzero product of the relevant functions
gives \(f\). Thus \(U=\pi^{-1}W\) is smooth affine. Apply the accepted
conditional smooth upper-growth gate to \(u|_U\), with the same
restricted witness. An affine FS exhaustion of \(U\) includes all
original source generators and the inverted functions. It dominates the
original upper bound. Conversely finite integrality bounds original
source generators by base generators: for a monic equation
\(z^r+\sum_{i<r}a_i z^i=0\), at every seminorm, \[
 |z|\le\max(1,\max_{i<r}|a_i|^{1/(r-i)}).
\] Polynomial degree bounds give the required rational affine upper
growth without changing either algebra.

Trace the resulting decreasing AFS sequence using \(T1\). This gives an
actual decreasing AFS sequence on \(W\), with limit \(N_\pi u\) there.
Positivity and \(\sum m_y=d\) give the upper bound \(Cd g+dB_0\) for its
limit. The accepted principal-open removal gate on the SMOOTH ORIGINAL
\(S\), retaining GEO, furnishes an actual affine FS limit \(v\) on \(S\)
agreeing with this trace on \(W\). No bad-fiber value has yet been
identified. Generic finiteness is preserved: on the generic good fiber
there is the generic trivial point, and its trace is finite.

\subsection{Slice at an arbitrary analytic
point}\label{slice-at-an-arbitrary-analytic-point}

Fix \(s\in S^{\rm an}\), set \(F=\mathcal H(s)\), and write its original
coordinates as \(\alpha\in F^n\). If \(s\in W^{\rm an}\) the conclusion
is already true. Otherwise choose \(z_0\in W(k)\) and set \[
 c(t)=\alpha+t(z_0-\alpha).
\] This is a nonconstant actual \(F\)-line. Work in the parameter open
\(T\subset\mathbb A_F^1\) where all base-chart denominators are nonzero;
it contains 0 and 1. The restrictions of \(a\) and \(f\) to the line are
nonzero polynomials, since their values at 1 are nonzero. Every
sufficiently small positive-radius parameter Gauss point belongs to
\(T\) and maps to \(W^{\rm an}\): a nonzero polynomial has nonzero Gauss
norm. The original image \(s_r\) tends to \(s\).

The original \(Y\) is geometrically integral because \(k\) is
algebraically closed. The line inverse image in \(Y_F\) is defined by
\(n-1\) affine linear equations. Krull's height bound gives height at
most \(n-1\) for every minimal prime; the finite-type domain dimension
formula over \(F\) gives component dimension at least one. Finiteness
over \(T\) gives dimension at most one. Thus, after reduction, EVERY
component is a horizontal curve. There are no isolated reduced vertical
components. Normalize this reduction to obtain \[
 q:C^\nu\longrightarrow T,\qquad p:C^\nu\longrightarrow Y_F.
\] Normalization is finite; each component is regular and smooth over
the perfect field \(F\). A nonconstant map to the regular parameter
curve is flat: its local modules over the base DVR are torsion free. The
total generic rank is exactly \(d\), since over the good open the entire
inverse image is the original finite étale base change and normalization
changes nothing there. For \(n=1\) there are no cutting equations and
the same argument applies. For \(n=0\), \(S,Y\) are single points and
the assertion is the direct degree-one identity.

External algebraic inputs, with every hypothesis used here, are
\href{https://stacks.math.columbia.edu/tag/0BXQ}{Stacks normalization
33.27.1}, \href{https://stacks.math.columbia.edu/tag/0BBZ}{height bound
10.60.12}, \href{https://stacks.math.columbia.edu/tag/02IJ}{dimension
formula 10.113.1},
\href{https://stacks.math.columbia.edu/tag/0B8Y}{normal curve smoothness
33.43.8}, and \href{https://stacks.math.columbia.edu/tag/0BXX}{curve
flatness 53.2.3}. The fields here have characteristic zero; finiteness
and the domain conditions are essential.

\subsection{Lift the entire curve to the ONE original
witness}\label{lift-the-entire-curve-to-the-one-original-witness}

Each normalized component maps generically into the isomorphism locus of
\(\mu_F\), hence has a rational lift. At every closed curve point the
local ring is a DVR. Apply the RELATIVE valuative criterion to
\(P_F\to Y_F\), using \(C^\nu\to Y_F\) as the base map. The unique local
extension spreads to a neighborhood by finite presentation.
Separatedness identifies overlaps. This gives a lift of the entire
\(C^\nu\) to the original \(P_F\), and it remains this lift after
further field extension. The target need not be proper over \(F\).
\href{https://stacks.math.columbia.edu/tag/0BX4}{Stacks 29.43.5} has
exactly this relative hypothesis; its pertinent proof uses 29.43.1 and
spreading of the local map.

At a lifted center restrict EACH original continuous contact approximant
before taking its decreasing limit. Its affine support becomes a real
linear combination of logarithms of analytic units, hence harmonic on
the smooth analytic curve. If the weight minus a harmonic test attains a
local maximum, its supporting harmonic function minus that test has the
same local maximum and is locally constant; the two inequalities force
the weight minus the test to be constant nearby. The harmonic-test
characterization and decreasing closure give subharmonicity, or the
identically-minus-infinity alternative on a local connected component.
The same applies to \(v\circ c\), using its actual decreasing AFS
approximants. No positive-current pullback or affine theorem over \(F\)
is invoked.

The curve inputs are exactly the pertinent units of
\href{https://theses.hal.science/tel-00010990v1/document}{Thuillier,
tel-00010990v1}: unit logarithms at LF3580--3627; harmonic maximum at
LF3730--3741; subharmonic comparison, closure and finite boundary values
at LF3783--3852; harmonic-test characterization at LF3860--3904. These
were retained by UTF8 byte decode and inclusive one-based LF split. A
local component containing a point with finite limit is not the
identically-minus-infinity alternative.

\subsection{An actual disc coordinate and a direct radial
limit}\label{an-actual-disc-coordinate-and-a-direct-radial-limit}

If \(F\) is trivially valued, first complete \(F(\tau)\) with an
independent nontrivial Gauss norm. In all cases pass isometrically to a
completed algebraic closure \(L\). All original pulled values stay
unchanged. The residue characteristic is zero: the trivially valued
subfield \(k\supset\mathbb Q\) embeds into the residue field. The finite
center fiber of the smooth \(C^\nu_L\) consists of \(L\)-rational
points.

Here is the analytic-coordinate argument rather than a formal-local-ring
substitution. At such a point smoothness gives a local algebraic
presentation \(F_i(z,w)=0\) with a square invertible \(w\)-Jacobian
\(J\) at the center, translated to 0. On sufficiently small closed
polydiscs the nonlinear part has \(w\)-Lipschitz norm strictly less than
one after multiplication by \(J^{-1}\); choose a smaller \(z\)-radius so
the constant \(w\)-term fits in that \(w\)-disc. The map
\(w\mapsto w-J^{-1}F(z,w)\) is then a strict contraction preserving the
disc. The successive iterates converge in its Tate algebra to an
analytic \(w(z)\). For every complete residue extension the same
contraction proves uniqueness of a solution. Thus the graph and
projection are inverse on the actual Berkovich neighborhood. A
presentation of the smooth local scheme is obtained by choosing a
nonvanishing Jacobian minor and shrinking its algebraic open; the
independent equations generate the local ideal there. This prevents
extra graph components. The identical contraction argument shows that
\(z+O(z^2)\) is an analytic coordinate after shrinking.

Let \(h\) be one of the pulled decreasing limits and \(x\) the center.
On a strict closed disc \(D_R\) in this coordinate its boundary is the
unique Gauss point \(\xi_R\). If \(h(x)\) is finite, the local limit is
subharmonic, its boundary value is finite by Thuillier 3.1.9, and
comparison with that constant gives \[
 h(x)\le h(\xi_R).
\] Upper semicontinuity and \(\xi_R\to x\) give the reverse limsup.
Hence \(h(\xi_R)\to h(x)\). If \(h(x)=-\infty\), usc gives this limit
directly: for every real \(M\) all sufficiently close values are below
\(M\). Local upper boundedness excludes a compensating plus-infinity
term. This proves a particular strict-disc radial limit, not just a
punctured-neighborhood limsup.

\subsection{Entire Gauss fibers, degrees, and canonical
regrouping}\label{entire-gauss-fibers-degrees-and-canonical-regrouping}

At each center \(x\in q^{-1}(0)(L)\) the map has form
\(t=b_xz^{e_x}U_x(z)\), where \(b_x\ne0\), \(e_x\ge1\), and
\(U_x(0)=1\). On a smaller disc \(|U_x-1|<1\). The binomial series for
an \(e_x\)-th root converges because all rational coefficients have norm
at most one and its argument has norm less than one. Replacing \(z\) by
\(zU_x^{1/e_x}\) is an analytic coordinate as just proved. The map is
EXACTLY \(t=b_xz^{e_x}\).

For small strict \(r\in|L^\times|\), the disc above \(\zeta_r\) has the
one Gauss point of radius \(R_x=(r/|b_x|)^{1/e_x}\), of completed
residue degree \(e_x\). To check uniqueness and degree, scale the radii
to one. Over the completed Gauss field the polynomial has residue
\(Z^{e_x}-T\), irreducible over \(\widetilde L(T)\) by Eisenstein at
\(T\). Its quotient is a field of degree \(e_x\), and a complete field
has a unique extension of its absolute value to a finite extension. Thus
there is one analytic point, with exactly that degree. Divisibility of
\(|L^\times|\) ensures \(R_x\) is also strict.

No missing points are concealed outside these coordinates. Fix a small
strict closed base disc in \(T_L^{\rm an}\). Its inverse image under the
finite \(q_L\) is compact: it is the spectrum of a finite algebra over
the corresponding disc affinoid algebra. Equivalently monic equations
bound finitely many algebra generators; all inverted chart denominators
are uniformly nonzero on the chosen base disc. In this fixed compact
inverse image take the closed complement of the union of ALL center
neighborhoods. Intersect it with inverse images of successively smaller
closed base discs. These compact sets decrease to the empty set, since
their intersection would lie over 0 and the full finite center fiber was
covered. Compactness makes one of them empty. Hence the center discs
account for the ENTIRE sufficiently small Gauss fiber, not only its
exhibited branches.

Before splitting the centers, the trace coefficient at a closed
\(F\)-point \(x\) of the normalized curve is \(e_x[\kappa(x):F]\); their
sum is \(d\). Why are these the ORIGINAL nonflat norm weights at 0? Pull
the original polynomial norm law to \(T\). Over its dense good open it
is the determinant norm of the finite étale cover. The product of
determinant norms of the finite flat normalized horizontal curves is a
regular polynomial law on \(T\), with the same law on that open.
Coefficientwise equality over an integral parameter ring gives equality
at 0. A uniformizer of order \(e_x\) makes the curve norm specialize to
its residue norm raised to \(e_x\). This identifies every residue
character, so grouping the centers by their images in \(Y_F\) gives
exactly \(m_y\). Subsequent complete-field compatibility groups split
\(L\)-centers back into each original analytic fiber point. No
fiber-dimension formula at a nonflat point is used. This also proves
independence from the chosen good line.

For all sufficiently small strict radii, therefore, \[
 v(s_r)=N_\pi u(s_r)
       =\sum_{x\in q^{-1}(0)(L)} e_x\,u\bigl(p(\xi_{x,r})\bigr).
\] The first equality was established on \(W\) BEFORE this argument.
Apply the radial limits to the left and every term. If all center values
are finite, ordinary finite summation applies. If one is minus infinity,
its positive-weight term tends to minus infinity and the other finitely
many terms remain bounded above, so the entire sum tends to minus
infinity. Regrouping yields \[
 v(s)=\sum_{y\mid s}m_y u(y)=N_\pi u(s).
\tag{T2}
\] This establishes the whole-base trace application under GEO on the
original ring and witness, including every nonflat, exceptional,
removed, and minus-infinity value. The norm-law part and smooth gate are
task-granted background with their binding restrictions, not invented
native support IDs. Singular EP, fiber descent, and the original
equation are not consequences of \(T2\).

Artifacts: \texttt{n1-trace.md}, combined \texttt{P16.1-argument.md},
standalone \texttt{P16.1-paper.tex} and \texttt{P16.1-paper.pdf};
primary units in \texttt{source-extracts/}, web retrievals in
\texttt{web-primary-01.json} through \texttt{web-primary-04.json},
premise coverage in \texttt{external-premises.json}. Reproduce with
\texttt{python3\ rerun.py}. All new combined deductions are unverified
drafts.

\section{n2. Controlled approximation, original-ring hulls, and a
countable
barrier}\label{n2.-controlled-approximation-original-ring-hulls-and-a-countable-barrier}

All claims in this note are unverified mathematical proof drafts. Except
in the explicitly countable barrier construction, retain an arbitrary
algebraically closed trivially valued field k of characteristic zero.
Y=Spec R is integral affine of finite type. An AFS function means a
rational finite maximum of positive multiples of original
regular-function logarithms and a rational constant floor, retaining a
positive-slope branch for EVERY chosen original ring generator. A
genuine AFS limit is an already supplied decreasing sequence of such
functions; membership is not inferred from singular EP. The
algebraic-first continuous contact/model/nef/FS bridge and genuine
decreasing closure are task-granted background in their exact round80
scope.

\subsection{G: a controlled original-reference
approximation}\label{g-a-controlled-original-reference-approximation}

Let g be an actual rational AFS exhaustion, shifted so g\textgreater=0,
and let v\_j decrease to a generically finite v with v\textless=g+B. Fix
rational c\textgreater1. We prove that v admits a decreasing
original-ring AFS sequence w\_j with a common upper bound cg+B', and
with \textbar w\_j-cg\textbar{} bounded on Y for each individual j.
Neither the bounds nor this conclusion assume singular envelopes.

First compactify the actual g. Clear its positive rational slopes with
an integer D. Use projective coordinates consisting of 1, the integer
powers occurring in Dg, AND all original generators z\_l. Take D large
enough that the extra branches
D\^{}(-1)log\textbar z\_l\textbar+lambda\_l are dominated by g;
polynomial degree bounds and the exhaustion bound allow rational
lambda\_l. The projective closure has the ORIGINAL affine chart Y, since
these unpowered generators generate R. Powers alone would not prove
this. Its boundary section s of H=O(1) and the actual rational bundle
B\_ref=H/D have an FS metric with affine weight exactly g. Write
g=g\_0+p, where g\_0 is the canonical boundary-section reference and p
the actual global FS potential, with its original cocycle.

Choose rational C\_i decreasing to minus infinity and rational R\_i
increasing to infinity, with (c-1)R\_i\textgreater B-C\_i+2. Let
Q\_i=\{g\textless=R\_i\}, which is compact. Choose increasing indices
j\_i such that v\_\{j\_i\}\textless g+B+1 on Q\_i. This follows because
the closed sets \{v\_j\textgreater=g+B+1\} in Q\_i decrease to the empty
set. It does NOT require Dini convergence to a discontinuous limit.

Define a\_i=max(v\_\{j\_i\},cg+C\_i) on Q\_i and a\_i=cg+C\_i outside
Q\_i. Near its entire outer boundary the second branch wins strictly, so
this is a continuous contact weight with a whole reference-winning
collar. The canonical potential is a\_i-cg\_0, not a\_i-cg; its boundary
value is C\_i+cp.~The granted bridge puts it in the genuine class
cB\_ref on the original compactification. These weights decrease: on
Q\_i both branches decrease; on Q\_\{i+1\}\textbackslash Q\_i the new
first branch is at most g+B+1\textless cg+C\_i, the old exterior value;
on the remaining exterior the floor decreases. Their limit is v on Y and
minus infinity relative to g\_0 at the boundary. Genuine decreasing
closure supplies membership before any sequential regularization is
invoked.

Genuine same-bundle FS regularization of this limit now gives w\_j. In
the section-s frame, compact boundedness of each continuous canonical
potential gives \textbar w\_j-cg\textbar=O(1); the first potential
bounds all later ones, giving the common upper bound. Alternatively,
uniformly approximate each continuous a\_i in the same bundle with
errors epsilon\_i=4\^{}(-i) and shifts b\_i=2
sum\_\{r\textgreater=i\}epsilon\_r. Since
b\_i-b\_\{i+1\}=2epsilon\_i\textgreater=epsilon\_i+epsilon\_\{i+1\},
these shifted approximants decrease. This also constructs the desired
sequence without EP.

Conversion to R is literal: a section on a tensor power divided by s to
that power is an original regular function. Append s\^{}m as the floor
and the sections representing z\_l times suitable powers of s at very
low rational weights. For sufficiently large common m they exist by
section localization. On the compact projective space their continuous
norms are bounded, whereas the approximating metric has a finite lower
bound; thus these added terms are dominated everywhere and change no
value. EVERY original generator remains, with positive slope 1/m. The
reference, tensor bundle and transition functions stay the same. The
construction works on a singular Y and a singular compactification for
an ALREADY genuine affine FS limit.

\subsection{Correct pole control through finite
normalization}\label{correct-pole-control-through-finite-normalization}

The printed denominator-order argument is false: in k{[}x{]}, f=x\^{}2
and p=x, p is not in (f) but vanishes on V(f). With q=3/2, b=p/f=1/x, q
log\textbar b\textbar\textless=max(0,log\textbar f\^{}(-1)\textbar),
while the printed q times denominator exponent is 3/2\textgreater1. The
reference G=max(0,log\textbar x\textbar) retains the original generator
and only enlarges this upper bound, so the counterexample satisfies the
genuine original-reference requirement. Its conclusion cannot be used.

Here is the replacement. Suppose b in R\_f is nonzero and q,c are
positive rational numbers, and

q log\textbar b\textbar{} \textless= max(G,c
log\textbar f\^{}(-1)\textbar)+B\_0

on Y\_f, where G is an original-ring AFS function finite everywhere on
Y. Choose N so m=Nc/q is an integer and put h=b\^{}N f\^{}m. Let Rbar be
the finite normalization of R. At each height-one prime E of Rbar
containing f, use the seminorm exp(-t ord\_E), t\textgreater0. It
restricts to an analytic point of Y\_f; every original regular function
has norm \textless=1, so G is bounded above independently of t. Divide
the displayed inequality by t and let t tend to infinity. It gives q
ord\_E(b)+c ord\_E(f)\textgreater=0, hence ord\_E(h)\textgreater=0. At
other height-one primes h is already regular, since h belongs to R\_f.~A
normal Noetherian domain is the intersection of its height-one
localizations, so h belongs to Rbar and is integral over R.

This uses precisely
\href{https://stacks.math.columbia.edu/tag/0BXQ}{finite normalization,
Stacks 33.27.1},
\href{https://stacks.math.columbia.edu/tag/031T}{height-one
intersection, 10.157.6}, and
\href{https://stacks.math.columbia.edu/tag/00PD}{the DVR
characterization}. R is a finite-type domain over a field, so its
normalization is finite and Noetherian; h is in its fraction field. All
hypotheses are retained.

The integral algebra R{[}h{]} is a finite R-module contained in
R\_f.~Choose a finite module generating list and clear its denominators
with ONE f\^{}M. Then

f\^{}M h\^{}j belongs to the ORIGINAL R for EVERY integer
j\textgreater=0.

No assertion that h itself, or a specified power of h, belongs to R is
needed. For example, R=k{[}t\textsuperscript{2,t}3{]}, f=t\^{}2,
b=t\^{}(-1), and h=bf=t is integral, with h\^{}2=f, but h is not in R.
The ONE multiplier f puts EVERY h\^{}j into R, since f
h\textsuperscript{j=t}(j+2) and every integer \textgreater=2 is in the
semigroup generated by 2 and 3. If x is in the original polynomial hull
of a compact L contained in Y\_f and x is also in Y\_f, apply its
defining inequalities to these original polynomials. They give

\textbar f(x)\textbar\^{}M \textbar h(x)\textbar\^{}j \textless= sup\_L
\textbar f\textbar\^{}M \textbar h\textbar\^{}j \textless= (sup\_L
\textbar f\textbar)\^{}M (sup\_L \textbar h\textbar)\^{}j.

If the last h-supremum is zero, the inequality already forces h(x)=0.
Otherwise take j-th roots and let j increase. Thus
\textbar h(x)\textbar\textless=sup\_L \textbar h\textbar. If
\textbar log\textbar f\textbar\textbar\textless=M\_1 on L union \{x\},
substitution of h gives

q log\textbar b(x)\textbar{} \textless= sup\_L q log\textbar b\textbar{}
+ 2cM\_1. (P)

The constant branch satisfies the hull bound directly. Apply (P)
branchwise to an AFS w\textless=max(g,c
log\textbar f\^{}(-1)\textbar)+O(1). Maxima commute with the finite
supremum, so the same 2cM\_1 error holds for w.

Now let v be an already genuine AFS limit on Y\_f, bounded above by an
original AFS reference on Y. Choose g large enough, including every
original generator, and set g\_c=max(g,c log\textbar f\^{}(-1)\textbar);
it is an actual reference for the ORIGINAL localized ring, including
f\^{}(-1). G with multiplier 2 supplies w\_i decreasing to v, with
w\_i\textless=2g\_c+B'. Formula (P) therefore gives
w\_i(x)\textless=sup\_L w\_i+4cM\_1. On a compact L, for decreasing
continuous w\_i and usc limit v, sup\_L w\_i decreases to sup\_L v.
Indeed, if a level exceeds the latter supremum, the nested closed
superlevel sets eventually become empty. This argument also covers
sup\_L v=-infinity. Let i increase and then c decrease to zero. We
obtain the exact original-ring hull inequality

v(x)\textless=sup\_L v. (H)

The case v(x)=-infinity is immediate. This proves principal-open hull
control, rather than a false denominator estimate.

\subsection{Countable entropy barrier, including every coordinate
face}\label{countable-entropy-barrier-including-every-coordinate-face}

Here assume k countable, so R is countable. Choose original generators
z\_1,\ldots,z\_k, enumerate every element f\_1,f\_2,\ldots{} of R, and
put g=max(0,log\textbar z\_l\textbar). For s coordinates t in (R union
\{+infinity\})\^{}s choose rational beta\textgreater=log(s+1), and
define

chi\_s(t)=beta\^{}(-1)log(1+sum\_i exp(-beta t\_i)), exp(-beta
infinity)=0.

Then max(0,-t\_i)\textless=chi\_s\textless=max(0,-t\_i)+1. At a point
with finite coordinate set I let p\_i=exp(-beta t\_i)/(1+sum exp(-beta
t\_l)), and p\_0=1/(1+sum exp(-beta t\_l)); omitted coordinates have
p\_i=0. Its support is ell=-sum\_\{i in I\}p\_i t\_i+C, with C chosen
for equality at that point. The entropy inequality proves
ell\textless=chi\_s. Equality is unique: on the finite face strict
log-sum-exp convexity uses p\_i\textgreater0 and p\_0\textgreater0; an
omitted coordinate becoming finite adds a positive term.

The support has a UNIFORM positive gap outside any neighborhood of the
point. On the finite face its recession expression is max(0,-a\_i)+sum
p\_i a\_i. It is positive for every nonzero vector a, since all p\_i and
p\_0 are positive. Hence the support difference is coercive there.
Restricting the finite coordinates to a compact box, any omitted
coordinate leaving a prescribed neighborhood of +infinity has a positive
lower bound on its exponential contribution. Compactness and strict
equality then give the claimed gap. At the all-infinity point the
support is zero and the same exponential argument applies. No
coefficient is placed on an infinite coordinate.

Here is global rational approximation, including those faces. On the
probability simplex in coordinates p\_0,\ldots,p\_s define
E(p)=beta\^{}(-1)sum p\_i log p\_i, with 0 log0=0. Triangulate it by
rational grids of mesh tending to zero, respecting ALL coordinate faces.
At each grid vertex choose rational h\_p in {[}E(p),E(p)+epsilon{]}. Its
convex roof is between E and E+delta, where delta tends to zero: the
lower bound is convexity, and the upper bound follows by interpolation
and uniform continuity on the compact simplex. Take its conjugate.
Explicitly

H\_r(x)=max\_\{grid vertices p\}(sum\_\{i\textgreater=1\}p\_i
x\_i-h\_p)+delta\_r

satisfies chi\_s(-x)\textless=H\_r(x)\textless=chi\_s(-x)+delta\_r
globally, with rational delta\_r tending to zero. All slopes are
nonnegative rational; all heights are rational. If x\_i=-infinity,
discard vertices with p\_i\textgreater0. The face triangulation gives
the SAME error bound, including the all-infinity face p\_0=1. This is
not an approximation proved only on a finite Euclidean box. Clearing
rational slopes converts each branch to an original monomial logarithm.
Rational majorants can be made decreasing by the summable-error shifts
used in G.

Let
rho\_j=chi\_\{k+j\}(-log\textbar z\_1\textbar,\ldots,-log\textbar z\_k\textbar,-log\textbar f\_1\textbar,\ldots,-log\textbar f\_j\textbar).
Polynomial degree bounds give rho\_j\textless=D\_j g+1, and
rho\_j\textgreater=g. Choose positive rational gamma\_j with
gamma\_j(D\_j+1)\textless=2\^{}(-j). Then

rho=sum\_\{j\textgreater=1\}gamma\_j rho\_j

converges uniformly on every compact generator-bounded set, is
continuous, nonnegative, and bounded by g+1. The upper partial sums
P\_r=sum\_\{j\textless=r\}gamma\_j rho\_j+2\^{}(-r)(g+1) decrease, since
the next summand is at most 2\^{}(-(r+1))(g+1). Uniform rational
approximations to their finite sums, with summable shifts, give a
decreasing actual AFS sequence for rho. Generator branches are already
dominated by rho\_j\textgreater=g and may be appended unchanged. This
proves membership directly from rational maxima and the granted
bridge/closure, not from the printed affine main theorem.

Given an analytic x and an open U containing x, countability and the
seminorm topology give a finite coordinate neighborhood with inverse
image contained in U. Choose j including those coordinates. Replace
rho\_j by its support ell at x and keep ALL the other positive summands.
The resulting rho' satisfies rho'\textless=rho, rho'(x)=rho(x), and
rho'\textless=rho-delta outside U, for some delta\textgreater0. Near x
its support contains only logarithms of units. Globally it can be
unbounded below or equal minus infinity; it has no plus-infinity values.

We must prove actual AFS membership of this replacement rather than
subtract a psh function. A positive real coefficient a times
log\textbar f\textbar{} is approximated by nested rational intervals
a\_r increasing to a and b\_r decreasing to a, with a\_r\textgreater0.
The functions max(a\_r log\textbar f\textbar,b\_r log\textbar f\textbar)
decrease to a log\textbar f\textbar: use b\_r for a positive logarithm,
a\_r for a negative one, and retain minus infinity at a zero. Zero
coefficients are omitted. Choose decreasing rational constants
approaching C. Approximate the other positive rho\_i-summands with the
same positive tail bounds and summable shifts. These explicit sums and
maxima are decreasing rational logarithmic expressions with the stated
pointwise limit, including minus infinity.

To retain a constant floor and EVERY original generator, append
max(epsilon\_r g-r, the expression), with positive rational epsilon\_r
decreasing to zero. The appended AFS expression decreases because
g\textgreater=0, and at every fixed analytic point epsilon\_r g-r tends
to minus infinity. It changes the LIMIT nowhere, including its
minus-infinity locus. It need not converge uniformly there, and uniform
convergence is not claimed. Restriction to a principal chart uses the
identical padding with an exhaustion of its actual localized ring,
including every new generator such as f\^{}(-1). Positive sums and
maxima of already genuine AFS limits use products of monomials and the
same padding. Thus rho' is constructed positively; subtracting rho-rho'
is never used as a positivity principle.

\subsection{Polynomial hulls and
Rossi}\label{polynomial-hulls-and-rossi}

For compact L the polynomial hull K=\{x:
\textbar f(x)\textbar\textless=sup\_L\textbar f\textbar{} for every f in
R\} is closed and compact: inequalities for the original generators
place it inside a compact affine polydisc domain. The AFS hull is the
same hull. One inclusion follows branchwise. For the other, a separating
original polynomial and rational threshold yield an AFS separator after
adding a very low constant floor and all original generator branches at
dominated low weights on L union \{x\}.

If compact polynomially convex K lies in an open U, first choose an
ambient generator-bounded compact domain with K in its strict interior.
Every point of its compact complement of U has a separating polynomial,
and a rational log threshold strictly between its value and its
K-supremum. Finitely many such strict separations, together with bounds
for every original generator, give a compact polynomial polytope P
contained in U and a STRICT polynomial open Q containing K, with Q
subset P. A strict finite polynomial domain is Runge: the suprema on any
compact subset are strictly below its defining thresholds, so its
polynomial hull stays in that strict domain. No equality between
arbitrary topological interior and the displayed strict inequalities is
used. Hulls respect closed immersions by the quotient-ring definition;
increasing unions of Runge opens are Runge by compactness.

Let z belong to hull(L)\textbackslash L and let V be relatively compact
open, with compact closure disjoint from L. Rossi asserts z belongs to
hull(K\_V), where K\_V=boundary(V) intersect hull(L). Suppose
log\textbar f(z)\textbar\textgreater t\textgreater sup\_\{K\_V\}log\textbar f\textbar.
Choose a nonnegative original AFS g\_0 equal to zero on a neighborhood
of hull(L), with compact zero-polytope in an open avoiding the bad part
of boundary(V) where log\textbar f\textbar\textgreater=t. It has a
strictly positive lower bound on that compact bad part. Choose rational
C so Cg\_0\textgreater log\textbar f\textbar-t on the ENTIRE boundary.
Glue max(Cg\_0,log\textbar f\textbar-t) on V to Cg\_0 outside. A whole
boundary collar is reference-winning. This continuous rational contact
weight equals the actual FS reference outside a compact set. On its
actual projective reference compactification the granted bridge gives
uniform same-bundle FS approximation, hence uniform original AFS
approximation after section conversion and generator retention. The
glued weight is zero on L and positive at z. A sufficiently accurate
approximant contradicts the AFS hull equality. This proves Rossi with
compact controlled closures; V itself is not called compact.

\subsection{The off-bad-set maximum
inequality}\label{the-off-bad-set-maximum-inequality}

Assume the ORIGINAL u is locally genuinely AFS regularizable on an
affine cover of Y\textbackslash Z, has original upper growth, and is
weak on one witness. It is not enough to know this only for u+c rho:
subtracting psh is unavailable. For rational c\textgreater0 put u\_c=u+c
rho. We prove u\_c(x)\textless=sup\_L u\_c for x in
hull(L)\textbackslash Z.

If this fails, take generators of the ideal of Z and a=max
log\textbar f\_i\textbar, with its empty-bad-set version a=0. Choose
small positive rational epsilon so u\_c+epsilon a still violates the
inequality at x. It is usc and upper bounded on compact hull(L), equals
minus infinity on Z, and has a finite attained maximum y outside Z union
L. Take V around y, with controlled compact closure in one good
principal affine chart, disjoint from L. Replace rho by rho' with
equality at y and gap delta on boundary(V). The explicit function w=u+c
rho'+epsilon a is an actual affine FS limit on that chart, using
ORIGINAL u's local regularizability and the positive constructions
above. Its upper growth comes from an original global AFS reference,
since rho'\textless=rho. It satisfies w(y)\textgreater sup\_\{K\_V\}w,
with gap at least c delta. Rossi puts y in the original hull of K\_V;
(H) on the principal chart gives w(y)\textless=sup\_\{K\_V\}w, a
contradiction. Every special value and the factor c are retained.

This note proves G, corrected original-ring pole control, the countable
barrier, Rossi, and the off-bad-set inequality. It does not assert
weak/global equality or singular EP. The task-granted background is
cited as an assumption, not as a registered native support. Diagnostic
Wu v1 LF911--1017,1603--1730,1815--1885,1965--2108 is retained to
identify the repaired dependencies; its printed theorems are not
premises. Artifacts: n2-affine-tools.md, sources/, source-extracts/,
external-premises.json. Rerun python3 rerun.py.

\section{n3. Conditional countable-field Runge, Hartogs, and
singular-stratum
chain}\label{n3.-conditional-countable-field-runge-hartogs-and-singular-stratum-chain}

This is a full mathematical proof draft of a CONDITIONAL countable-field
chain, using n1 and n2. The field k is countable, algebraically closed,
trivially valued, and of characteristic zero. Y is integral affine of
finite type. The original u is usc, generically finite, locally bounded
above, has an upper bound by an actual original-ring AFS exhaustion, and
is weak on ONE proper birational contact witness. All uses of the
accepted smooth affine gate retain GEO as defined in n1, including the
required smooth ample-boundary compactification and projective
birational domination of that same restricted witness by a smooth
projective source. No arbitrary completed-residue-field affine theorem
or singular EP is imported.

\subsection{Split pieces on one controlled neighborhood, including
gamma=1}\label{split-pieces-on-one-controlled-neighborhood-including-gamma1}

Suppose pi:Y-\textgreater S is finite surjective onto a principal affine
open of A\^{}n.~Suppose Z is closed, pi\textbar Z is a CLOSED IMMERSION
onto its image, and Z and Z'=pi\^{}(-1)(pi(Z))\textbackslash Z are
disjoint closed subsets. Assume the ORIGINAL u is locally genuinely AFS
regularizable off Z. We prove its strict sublevels are Runge.

Use n2's countable nonnegative rho and its positive support
replacements. If a strict sublevel of u\_c=u+c rho, c\textgreater0
rational, is not Runge, there is compact L such that its compact
polynomial hull K contains a point where u\_c\textgreater=0 although
sup\_L u\_c\textless0. After an additive normalization let sup\_K u\_c=0
and sup\_L u\_c=-delta\textless0; if the latter is minus infinity choose
any finite delta with sup\_L u\_c\textless=-delta. The finite attained
maximum x belongs to Z by n2's off-bad-set inequality.

Put C=hull\_S\{pi(x)\}. Because pi\textbar Z is a closed immersion, the
quotient-ring hull property gives C subset pi(Z)\^{}an. The ENTIRE
compact preimage of C consists of disjoint compact pieces C\_1 in Z and
C\_2 in Z'. Choose disjoint open neighborhoods A\_1,A\_2 of these
pieces, with A\_1 avoiding Z' and A\_2 avoiding Z. Properness of the
finite analytic map gives an open neighborhood A of C whose entire
preimage is in A\_1 union A\_2: take the complement of the closed image
of Y\^{}an(A\_1 union A\_2).

By n2's strict polynomial neighborhoods choose a nonnegative actual base
AFS g with C in an inner strict neighborhood U where g=0, and
\{g\textless=1\} compactly contained in A. Let W=\{g\textless1\}; its
preimage splits into the clopen OPEN pieces W\_i=pi\^{}(-1)(W) intersect
A\_i. These interiors are not compact; their controlled closures in
pi\^{}(-1)\{g\textless=1\} are compact. On the preimage of W put h=1 on
W\_1 and 0 on W\_2. The global continuous contact weight H=max(3g circ
pi-1,h) there, and 3g circ pi-1 elsewhere, has a whole reference-winning
outer collar, and equals h over U.

The task-granted continuous canonical norm trace makes the trace of this
indicator over U continuous and integer valued:

gamma\_1(s)=sum\_\{y in pi\^{}(-1)(s) intersect W\_1\}m\_y.

It is only LOCALLY constant. Shrink ONE base neighborhood N of pi(x)
inside U where gamma\_1 is the positive integer gamma; retain this N for
EVERY subsequent hull, boundary and fiber test. All weights m\_y are the
raw positive canonical integers of n1, never fiber lengths. A closed
immersion guarantees analytic fiber injectivity, so the W\_1 fiber above
pi(x) is exactly \{x\}, even after a completed residue-field extension.

For gamma\textgreater1 shrink N further so all its W\_1 fiber values of
u\_c are \textless delta/{[}2(gamma-1){]}. This follows from usc at the
unique central point x and properness: cover that point by the requisite
strict sublevel, cover the other central piece by W\_2, and exclude the
closed image of the complement. For gamma=1 do not divide by gamma-1.
Positivity and integrality imply every W\_1 fiber over N is one point of
weight one.

Choose a relatively compact open V around x with compact closure in
pi\^{}(-1)(N) intersect W\_1, disjoint from L. Let K\_V=boundary(V)
intersect K. Replace rho by rho' with equality at x and
rho'\textless=rho-delta' on boundary(V). Put u'\_c=u+c rho'. Then
u'\_c(x)=0, u'\emph{c\textless=u\_c, and sup}\{K\_V\}u'\_c\textless=-c
delta'. The factor c is essential.

On W\_1 put v=u'\_c, on W\_2 put v=0. On the controlled compact closure
these have a finite common upper bound M. Choose a rational
C\_0\textgreater M+1+delta and define

tilde v=max(C\_0 g circ pi-delta,v) on pi\^{}(-1)(W), tilde v=C\_0 g
circ pi-delta outside.

For g sufficiently close to 1 the reference wins on the ENTIRE collar.
Pullbacks, products, support replacements and maxima all use the SAME
original witness. Its pullback remains locally a decreasing contact
limit. The global upper bound is an actual original reference, and the
floor makes tilde v finite. Over U it is max(u'\_c,-delta) on W\_1 and
zero on W\_2. Apply n1, retaining its GEO applications, to obtain
w=N\_pi(tilde v) as a genuine original-base AFS limit. Pullback w circ
pi is a genuine original-source AFS limit. To retain every source
generator, use its monic norm bound to dominate a source exhaustion by a
positive multiple of a base exhaustion and append decreasing low
floor/generator branches as in n2; the limit is unchanged at every
point.

Rossi gives x in hull\_Y(K\_V), hence w(pi(x))\textless=sup\_\{K\_V\}w
circ pi. Its left side is zero. For z in K\_V intersect Z, analytic
injectivity and avoidance of Z' make the W\_1 fiber exactly \{z\}, of
weight gamma, giving

w(pi(z))\textless=-gamma min(delta,c delta').

For y in K\_V\textbackslash Z, the off-bad-set inequality gives
u\_c(y)\textless=sup\_L u\_c\textless=-delta, so tilde v(y)=-delta.
Every other W\_1 fiber point is in the SAME pi\^{}(-1)(N); when
gamma\textgreater1 its value is at most delta/{[}2(gamma-1){]}. Thus,
writing its positive weight m\_y\textgreater=1,

w(pi(y)) \textless= -m\_y delta+(gamma-m\_y)delta/{[}2(gamma-1){]}
\textless= -delta/2.

When gamma=1 there is no other W\_1 point and w(pi(y))=-delta. The W\_2
values are zero. These are uniform negative bounds on ALL of K\_V, a
contradiction. Consequently each \{u+c rho\textless0\} is Runge. As
rational c decreases to zero these opens increase to \{u\textless0\},
since rho is finite and nonnegative; compactness gives Runge for their
union. The argument applies at every translated level. It preserves
minus infinity, all raw weights, the whole collar, and one controlled
base neighborhood.

\subsection{Genuine Runge exhaustion before coefficient
radii}\label{genuine-runge-exhaustion-before-coefficient-radii}

For a countable ring R, the seminorm topology of Y\^{}an is a subspace
topology of the countable product of the norm coordinates
\textbar f\textbar, f in R. It is second countable and metrizable.
Generator-bounded polydisc domains are compact and exhaust Y\^{}an
through their strict interiors. Hence every open Runge D admits a
countable exhaustion by compact subsets L\_r with controlled interiors.

Inductively take the polynomial hull of L\_r union the previous
controlled closure. It is compact and stays in D by the Runge property.
Apply n2's thick polynomial-neighborhood construction to obtain STRICT
finite polynomial opens Q\_r, with compact controlled closures P\_r
contained in D, satisfying P\_\{r-1\} subset Q\_r and L\_r subset Q\_r.
Then Q\_r increase and exhaust D. Their defining polynomials retain
bounds for EVERY original generator. We use
Q\_r=\{log\textbar P\_\{rj\}\textbar\textless b\_\{rj\}\}, with rational
b\_\{rj\}; we do not identify the interior of an arbitrary closed
weak-inequality domain with its strict-inequality locus.

Apply this to a genuinely established Runge D in (Y times
A\textsuperscript{1)}an containing the zero section, writing each
P\_\{rj\}=sum\_i a\_\{rji\} w\^{}i in ORIGINAL R{[}w{]}. The opens where
zero lies in Q\_r are

Omega\_r=\{log\textbar a\_\{rj0\}\textbar\textless b\_\{rj\} for all
j\};

they increase and cover Y\^{}an. For y in Omega\_r the fiber is the
analytic affine line over F=H(y). Define its open-disc containment
radius d\_r(y). The EXACT formula is

-log
d\_r(y)=max\_\{j,i\textgreater=1\}(log\textbar a\_\{rji\}(y)\textbar-b\_\{rj\})/i.
(D)

Zero coefficients contribute minus infinity; an empty maximum is minus
infinity and means radius infinity. The constant-coefficient conditions
defining Omega\_r are indispensable.

To prove (D), on a closed radius-s disc the polynomial Gauss norm is
max\_i \textbar a\_i\textbar s\^{}i at EVERY positive real s, whether or
not s is in the value group. The ultrametric inequality bounds every
point by that norm. Hence the open radius-r disc is contained in all
strict constraints exactly when each positive-degree coefficient has
\textbar a\_i\textbar r\^{}i\textless=exp(b\_j), given the strict
constant constraints. Equality at the endpoint causes no problem, since
the endpoint is excluded. Conversely a larger open radius contains a
Gauss point violating a constraint. This proves (D) over any completed
residue field using only polynomial norms, not a regularization theorem
there.

By the nesting, d\_r increase. For every s\textless d\_D(y), the compact
closed fiber disc of radius s lies in D, so a finite subcover of the
increasing Q\_r puts it in one Q\_r. Thus d\_r(y) increase to d\_D(y),
including infinity, and A\_r=-log d\_r decrease to -log d\_D on every
eventually defined Omega\_r. The exhaustion includes a w-bound, so
individual A\_r can be taken to have a constant floor. Their rational
positive coefficient maxima are continuous contact weights on Omega\_r.

If u\_D=-log d\_D is generically finite and bounded above by an original
AFS reference g+B, we now prove controlled ORIGINAL-ring regularization.
Choose R\_i increasing and C\_i decreasing as in n2's G, and indices
r\_i increasing so the compact \{g\textless=R\_i\} lies in
Omega\_\{r\_i\} and A\_\{r\_i\}\textless g+B+1 there. The latter
assertion follows from nested compact closed superlevel emptiness; no
Dini assertion for u\_D is used. Glue max(A\_\{r\_i\},cg+C\_i) inside
\{g\textless=R\_i\} to cg+C\_i outside. The outer reference wins on a
whole collar. These continuous contact weights decrease on the old
compact domain, the new annulus and the exterior separately, exactly as
in G. The granted bridge and decreasing genuine closure give membership
in ONE actual bundle cB\_ref, then same-class FS regularization gives
original-ring affine FS weights decreasing to u\_D. Every original
coefficient and every original generator is retained by localization and
dominated padding. Infinite radius values give minus infinity correctly.
This proves the needed Hartogs conclusion after Runge, without
presupposing the affine weak/global theorem.

\subsection{Hartogs on the same product
witness}\label{hartogs-on-the-same-product-witness}

For any original u as above put H=log\textbar w\textbar+u circ p on Y
times A\^{}1, defining its value to be minus infinity when w=0 or
u=-infinity. There are no plus-infinity values. Its one witness is P
times A\^{}1. Restrict each original contact approximant and approximate
log\textbar w\textbar{} by its decreasing rational floors on the same
neighborhoods before taking sums and limits. Supports away from its
zeros are unit logarithms, and the zero floors provide continuous
contact weights. Thus H is weak on that SAME product witness. It has
upper growth by a product exhaustion retaining every Y-generator and w.

On a good affine chart where ORIGINAL u is an AFS limit, H is an actual
AFS limit: form products for positive sums, use decreasing floors for
log\textbar w\textbar, and append decreasing low generator branches
retaining every product-chart generator. If the split-piece hypotheses
hold on Y, they hold on Y times A\^{}1 with Z times A\^{}1 and pi times
identity, including closed immersion and closed disjoint complementary
piece. Applying the already proved split argument to H proves

D=\{log\textbar w\textbar+u circ p\textless0\}

is genuinely Runge in the PRODUCT. This is not inferred merely from
sublevels on Y. It contains the zero section, and its fiber radius is
exp(-u(y)), with infinity when u(y)=-infinity. Therefore the preceding
coefficient argument gives actual AFS regularization of u. The same
reasoning with EMPTY bad set will supply the final local-to-global step.

\subsection{Descending Z and
localization}\label{descending-z-and-localization}

Suppose only pi\textbar Z is a closed immersion; drop the initial
disjoint-piece hypothesis. Define reduced closed sets Z\_1=Z and

Z\_\{r+1\}=Z\_r intersect
closure\_Y(pi\^{}(-1)(pi(Z\_r))\textbackslash Z\_r).

They strictly drop dimension until empty. To check ALL components, a
closed immersion identifies pi(Z\_r) with Z\_r; images of its
irreducible components are distinct and incomparable. Each component T
of Z\_r is a component of pi\^{}(-1)(pi(Z\_r)): a containing component
maps onto a containing component of pi(Z\_r), and finiteness preserves
its dimension, forcing equality. The closure of the complementary union
therefore omits the generic point of every component T. Its intersection
with Z\_r has strictly smaller maximum dimension. Outside
pi\^{}(-1)(pi(Z\_\{r+1\})), the complementary piece
pi\^{}(-1)(pi(Z\_r))\textbackslash Z\_r is closed, since the only new
closure points on Z\_r are in Z\_\{r+1\}. This also proves the required
localization equality, rather than just a set-theoretic dimension
estimate.

Inductively let Pi\_r assert ORIGINAL local AFS regularizability at
every closed point of Y\textbackslash Z\_r. Pi\_1 is the assumed off-Z
cover. A closed y in pi\^{}(-1)(pi(Z\_\{r+1\}))\textbackslash Z\_\{r+1\}
is off ORIGINAL Z: otherwise closed-immersion injectivity places it in
Z\_\{r+1\}. Use the original off-Z cover. For any other y in
Y\textbackslash Z\_\{r+1\}, choose a principal affine base S' inside
S\textbackslash pi(Z\_\{r+1\}) around pi(y), and set U=pi\^{}(-1)(S')
EXACTLY. The induced finite map is onto S'. On U the bad set U intersect
Z\_r and its complementary inverse-image piece are disjoint closed; the
off-bad-set local hypothesis is Pi\_r. Apply the split proof on U times
A\^{}1 using the same product witness. Its Hartogs domain is Runge, so
the preceding coefficient argument gives actual AFS regularization on U.
This proves Pi\_\{r+1\}.

At the empty last Z\_r we have an algebraic affine local cover.
Closed-point neighborhoods cover every scheme point: finite-type schemes
over a field are Jacobson, so a nonempty closed complement would contain
a closed point (\href{https://stacks.math.columbia.edu/tag/0478}{Stacks
33.14.1}). Their analytifications therefore cover every analytic point.
A local cover is not yet the desired global affine conclusion. Apply
n2's off-bad-set proof with Z empty to get every \{u+c rho\textless0\}
Runge, hence \{u\textless0\} Runge. Apply that SAME empty-bad-set
reasoning to H on the product; now its Hartogs domain is Runge, and the
coefficient argument gives ORIGINAL global AFS regularization of u. This
supplies the final LOCAL-TO-GLOBAL implication that was missing from the
printed induction.

\subsection{Finite projection along a smooth stratum, with the correct
U}\label{finite-projection-along-a-smooth-stratum-with-the-correct-u}

Let Y be integral affine of dimension n and Z proper smooth closed, with
dim Z\textless=n-1. For a closed k-point z in Z, embed the projective
closure X into P\^{}D with D\textgreater=n+1. Let B=X\textbackslash Y
and T=closure\_X Z union B, reduced. T has dimension \textless=n-1 and
is smooth at z, since z is outside B.

Choose a projection center p outside X, outside the closed secant cone
from z to T\{z\}, and outside the projective tangent space of T at z.
The secant cone has dimension \textless=dim T+1\textless=n\textless D.
The tangent space and X also have dimension \textless D, so a k-point
outside their union exists. Projection from p is finite on X by
\href{https://stacks.math.columbia.edu/tag/0B1P}{Stacks 43.23.1}. Its
T-fiber over the image of z is only z, and its differential on T there
is injective. This is the secant/tangent argument of
\href{https://stacks.math.columbia.edu/tag/0B1Q}{43.23.2}; the following
local algebra proves the needed statement also for the reduced union T
with possibly several components away from z.

Let C be the local ring of the scheme-theoretic image at q, and D\_T the
finite local ring above it. There is just the one point above q, with
residue field k. Injective differential gives a surjection on cotangent
spaces, so m\_\{D\_T\}=m\_C D\_T+m\_\{D\_T\}\^{}2. Nakayama gives
m\_\{D\_T\}=m\_C D\_T. Thus D\_T/m\_C D\_T=k is generated by 1; Nakayama
for the finite C-module D\_T gives D\_T=C. The image inclusion and
projection are consequently inverse over a Zariski neighborhood of q,
after shrinking to kill the finite kernel and cokernel. Here T is smooth
at z, and the secant exclusion gives the single fiber; it is not enough
to know an arbitrary set-theoretic injectivity statement.

Repeat the projection while the ambient dimension exceeds n.~At each
step the image T is still smooth at the tested point through the just
proved local isomorphism, and has dimension \textless=n-1; the integral
projective image of X has dimension n.~The final image of X in P\^{}n is
all P\^{}n.~The tested point is outside the image of B, since B is part
of T and the T-fiber was the single z. Choose a principal affine open S
around its image, avoiding the image of B and lying in the isomorphism
neighborhood of T. Set

U=r\^{}(-1)(S), pi=r\textbar U:U-\textgreater S.

This is the correct U, not the printed r\^{}(-1)(W) with a different
target S. U is affine, is contained in the ORIGINAL Y, pi is finite onto
S, and pi\textbar Z\_U is an isomorphism onto a closed image. Its
analytic restriction is injective because a closed immersion stays a
closed immersion after every complete field extension. Restricting u and
its original witness preserves the hypotheses; monic integrality bounds
give the required reference growth on S. This geometric
finite-projection step is proved here. It does not discharge GEO's
separate compactification/resolution antecedents.

\subsection{Singular-stratum induction and the conditional affine
conclusion}\label{singular-stratum-induction-and-the-conditional-affine-conclusion}

Let Y\_0=Y and Y\_\{r+1\} be the reduced singular locus of Y\_r. In
characteristic zero the regular locus is smooth and dense in EVERY
reduced irreducible component
(\href{https://stacks.math.columbia.edu/tag/0B8X}{Stacks 33.25.8}), so
maximum dimension strictly drops until Y\_r is empty. Reducible and
unequal-dimensional strata are allowed. Let Sigma\_r assert local
genuine AFS regularizability at all closed points of
Y\textbackslash Y\_r. Sigma\_1 follows on smooth affine opens from the
accepted smooth gate, retaining all its GEO hypotheses and restrictions
of the SAME original witness.

At a new point y in Y\_r\textbackslash Y\_\{r+1\}, choose an affine
neighborhood V inside Y\textbackslash Y\_\{r+1\}. Its Z=V intersect Y\_r
is a smooth proper closed stratum. The finite projection just proved
supplies U around y and pi:U-\textgreater S with pi\textbar Z\_U a
closed immersion. The induction gives ORIGINAL local AFS
regularizability on U\textbackslash Z\_U. The
descending-Z/product-Hartogs theorem supplies actual AFS
regularizability on U. This proves Sigma\_\{r+1\}. At the final empty
stratum, invoke the EMPTY-Z local-to-global conclusion above, not merely
the existence of a local cover.

Consequently, under GEO for its specified smooth applications, every
original upper-growth weak input on integral affine Y over countable k
is an ORIGINAL-ring decreasing AFS limit, and every strict sublevel is
Runge. Every whole-base trace used in this chain is n1's trace on the
one original witness. No singular envelope theorem, Wu affine main
theorem, arbitrary-field affine regularization, or printed
denominator-order claim is a premise.

\subsection{A limited countable global membership
consequence}\label{a-limited-countable-global-membership-consequence}

Let X be integral projective over countable k, L ample rational, and phi
generically finite, globally upper bounded and weak relative to the
canonical reference of L. Require the same contact/GEO hypotheses on
each of its specified affine witnesses; a common domination, if used,
remains a conditional geometric antecedent. Choose a very ample tensor
qL with homogeneous sections z\_i. Its principal charts Y\_i have actual
reference weights rho\_i=q\^{}(-1)max\_j
log\textbar z\_j/z\_i\textbar\textgreater=0. Then u\_i=phi+rho\_i has
upper growth and the conditional affine theorem applies.

Use G on the ORIGINAL projective compactification X, with its actual
section z\_i and matching reference rho\_i. For each rational
c\textgreater1 this proves genuine membership in cL of the function
phi-(c-1)rho\_i on Y\_i, minus infinity outside. At every analytic point
one of the homogeneous coordinates attains the maximum; there rho\_i=0.
Therefore the maximum over all chart functions equals ORIGINAL phi on
all of X\^{}an, including minus infinity. Finite maxima and
varying-class decreasing closure as c decreases to one prove phi in
PSH(L).

This is conditional countable weak/global membership for the stated
actual ample canonical reference. It does NOT show that every genuine
resolved-nef potential is constant on the analytic fibers of a
resolution, and so it supplies no singular EP by itself. Onto birational
descent cannot be inferred from EP on the resolution: BJ5.13 needs
approaching EP on the TARGET.

Evidence and reproduction: n1-trace.md and n2-affine-tools.md contain
the required trace and elementary tools. Diagnostic Wu v1
LF1731--1814,1886--1964,2200--2696 and its envelope proof LF2697--2755
are pinned; their printed conclusions are not granted premises. The
primary projection units are retained in web-primary-03.json and
external-premises.json. Run python3 rerun.py for the complete paper,
both TeX passes, hashes and finite documentary checks. All conclusions
remain unverified drafts; the unchanged original P16.1 is not solved.

\section{n4. Exact equation domain, source audit, and the remaining
gates}\label{n4.-exact-equation-domain-source-audit-and-the-remaining-gates}

Candidate=false. The unchanged original P16.1 and its entire
P16.2--P16.7 family remain exactly those preserved in task-original.json
and input-original.json. The new countable conditional affine theorem
does not replace the original arbitrary-cardinality
singular/geometrically unibranch projective equation domain.

The retained equation domain is reduced irreducible geometrically
unibranch projective X over an ARBITRARY algebraically closed trivially
valued characteristic-zero field k, an ACTUAL ample rational line bundle
L, every dimension, and positive Radon targets with arbitrary support
and the intrinsic finite-dual-energy condition. Nothing here establishes
arbitrary nontrivially valued fields, positive residue characteristic,
nonample classes, proper nonprojective positivity, or nonalgebraizable
spaces. The special contact class and GEO in n1--n3 remain binding
assumptions of those intermediate results.

\subsection{What the new drafts establish and what is
conditional}\label{what-the-new-drafts-establish-and-what-is-conditional}

n1 reconstructs the whole-base trace-value application at every original
analytic point under the accepted smooth gate's retained GEO
antecedents. The base extension is constructed first on the original
smooth normal base. Every normalized horizontal slice component, every
residue degree and every canonical norm character is retained. Relative
properness lifts the curves to the ONE original witness. Actual
contraction coordinates and strict-disc radial limits include minus
infinity; compactness covers the ENTIRE finite Gauss fiber. Original
nonflat weights are recovered by equality of polynomial norm laws and
complete-field regrouping, never by fiber dimensions.

n2 proves controlled approximation G for an ALREADY genuine AFS limit on
a possibly singular variety, with original reference, whole collars,
original generators, and no singular EP. Its finite-normalization pole
repair replaces the false denominator-order claim by height-one bounds
and ONE uniform f\^{}M putting every power back in the ORIGINAL ring.
Its entropy-grid approximation treats all coordinate faces, rational
vertex heights/slopes, global errors, rational coefficient intervals,
countable tails and summable shifts. Rossi and off-bad-set control use
the ORIGINAL u's local regularizability; no subtraction of psh is used.

n3 gives the conditional countable-field split-piece proof with only
locally constant integer branch sums, one shrunken controlled base
neighborhood, raw positive weights, the factor c in the boundary gap,
and a separate gamma=1 case. It proves descending Z, localization,
genuine Runge exhaustion BEFORE strict coefficient radii, infinite
radius values, original coefficient regularization, finite projection
along a smooth stratum with EXACT U=pi\^{}(-1)(S), singular-stratum
induction, and the final EMPTY-Z local-to-global step. Products and
maxima retain the original witness. The finite-projection geometry is
reconstructed; GEO's other
compactification/projective-domination/resolution applications are NOT
independently discharged. A canonical ample-reference countable
weak/global membership consequence is correspondingly conditional.

These are mathematical proof drafts, not independently reviewed results.
No root86/root87 draft, printed Wu affine main theorem, printed Wu
envelope conclusion, or arbitrary completed-residue-field regularization
theorem was promoted to a premise. Task-granted round77/80 bodies are
background assumptions, not registered native support IDs and not
negative-index external premises.

\subsection{Genuine global operator and exact raw equation
conventions}\label{genuine-global-operator-and-exact-raw-equation-conventions}

The task-granted global source audit of Boucksom--Jonsson, 1801.08229v4,
is confirmed against its exact authorized primary text. Its convention
permits reduced possibly nonnormal projective varieties; Section 7
resumes general projective scope. Relative integral closure of a test
configuration is INSIDE its prescribed generic fiber, preserving that
original nonnormal fiber. Its raw model measure is

\[
 \operatorname{MA}^{\rm raw}(\phi)
 =\sum_E b_E\bigl((\omega_{\mathcal X}+D)|_E\bigr)^n\delta_{v_E},
 \qquad b_E=\operatorname{ord}_E(t),\quad
 v_E=b_E^{-1}\operatorname{ord}_E|_{k(X)}.
\]

Central fundamental-cycle multiplicities, vertical DVRs, and the actual
bundle are essential. The genuine mixed finite-energy pairing constructs
positive Radon measures of raw mass V=(L\^{}n), with genuine decreasing
continuity. This is an existing genuine operator on general reduced
projective varieties, not an inverse Banach-algebra measure
construction.

Use the normalized energy

\[
 E_\omega(\phi)=\frac{\langle(\omega,\phi)^{n+1}\rangle}{(n+1)V},
 \qquad\mathcal E^1(\omega)=\{\phi\in\operatorname{PSH}(\omega):E_\omega(\phi)>-\infty\}.
\]

A raw target nu must be positive Radon with nu(X\^{}an)=V. For mu=nu/V
its precise energy requirement is

\[
 E_\omega^\vee(\mu)=\sup_{\psi\in\mathcal E^1}
       \left(E_\omega(\psi)-\int\psi\,d\mu\right)<\infty.
\]

BJ9.20 equates this with integrability of EVERY genuine finite-energy
potential on irreducible X. There is no support restriction. BJ10.4 and
10.6 give uniqueness modulo one real constant and genuine domination
without assuming EP. In dimension zero X is one point, V=1, and the
empty product and constant ambiguity have their literal meaning.

The actual general-variety theorem is
\href{https://arxiv.org/pdf/1801.08229v4\#page=121}{BJ12.8}: onto
finite-energy Monge--Ampere is equivalent to EP. Under EP, its
variational proof uses normalized weak PSH compactness, a bounded-energy
estimate, finite-energy target integral continuity (9.19), and upper
semicontinuity of energy to obtain an ACTUAL maximizer. The
differentiated envelope energy identifies its genuine measure. The raw
functional VE\_omega(phi)-integral phi dnu is constant invariant exactly
at mass V. This proof cannot be replaced by an I-Cauchy sequence, an
abstract completion, or approximate maximizers without attainment.
Smooth EP on a resolution alone is not target EP.

Exact primary units retained include BJ LF435--469 (variety convention),
LF3649--3692 (general scope), LF3980--4148 (operator/mass/decreasing
continuity), LF5230--5304 (integral continuity and intrinsic energy),
LF5440--5479 (uniqueness/domination), LF6092--6146 (actual EP/onto
theorem and proof), LF6230--6268 (sequential regularization AFTER
membership), and LF2859--2876 (target-side birational descent).
Coordinates are inclusive one-based UTF8 bytes decode
split(`\textbackslash n'). The catalogue manifest identifies any broader
retained range containing these units; source files and literal bodies
are preserved. These citations describe task-granted background, not a
new singular existence assertion.

\subsection{Concrete indispensable
gaps}\label{concrete-indispensable-gaps}

\begin{enumerate}
\def\labelenumi{\arabic{enumi}.}
\item
  GEO is still conditional in its stated applications: smooth
  ample-boundary compactifications and projective birational
  domination/resolution of the ONE original witness with matching actual
  references have not been proved here from exact primary geometric
  statements. The analytic deductions retain them rather than silently
  discharging them.
\item
  Genuine resolved-nef fiber-value descent has not been established. It
  requires constancy of every relevant genuine pulled-nef potential on
  all geometrically connected analytic fibers after EVERY completed
  residue-field extension, including minus infinity; actual
  varying-class references, restrictions through curve MORPHISMS, and
  minimum/usc descent must be checked. The affine theorem over countable
  k cannot be reapplied over arbitrary H(s). BJ5.13 cannot supply onto
  descent from EP on the smooth source: it requires approaching EP on
  the TARGET.
\item
  Singular countable EP is therefore not proved.
  \href{https://arxiv.org/pdf/2609.27610v1\#page=51}{Wu2609.27610v1's
  envelope proof}, retained at LF2697--2755, explicitly calls its
  resolved-nef Proposition2.24 and weak/global Theorem6.5. n3 repairs a
  conditional version of the latter gate; it does not independently
  validate the former or the resulting envelope.
\item
  The original ARBITRARY-CARDINALITY field domain still requires
  finite-data obstacle/competitor descent and ascent. Relative integral
  closure must preserve the original nonnormal generic fiber, vertical
  R1/S2 behavior, component multiplicities, positive raw intersection
  masses, reference compatibility, orthogonality and domination.
  Different finite competitors can require different larger countable
  fields. An arbitrary collection of PSH competitors does not
  simultaneously descend merely because one finite model does.
\item
  Normalization/geometrically unibranch descent still requires
  approaching target-side EP, degree-one fundamental-cycle pairing
  projection, and compatibility of energies and the genuine E1/MA
  operator through unbounded limits. It is not enough to identify
  analytic point sets or continuous model measures.
\item
  Only after these gates can the actual singular envelope
  compactness/differentiation/attainment proof yield MA\_raw(phi)=nu for
  EVERY prescribed finite-dual-energy raw mass-V target. The operator
  and uniqueness already available in general scope do not provide this
  missing existence assertion.
\end{enumerate}

\subsection{Documentary execution, not correctness
certification}\label{documentary-execution-not-correctness-certification}

The exact input bytes are preserved. The 23 authorized primary files and
54 exact LF units were hash checked and copied; public primary Stacks
statement/proof snapshots are additionally retained with their own
hashes, actual URLs and nonnegative premise indices. Web snapshot LF
means the stored UTF8 tool-rendered text's LF coordinates, not imaginary
source-HTML coordinates or an immutable Stacks edition. Cached arXiv/HAL
versions remain exact. External-premises.json covers the external
geometric/algebraic/radial facts used; task-granted full bodies are
recorded separately as background.

Run python3 rerun.py in this workspace. It checks preservation and
source/unit hashes, regenerates the complete argument and native
notes/status, runs finite rational diagnostics, invokes the EXISTING
Pandoc and pdfTeX twice, and saves commands, both pass logs, outputs,
environment and artifact hashes. finite-checks.json declares exactly
which finite cases were tested. These are finite observations only. No
compilation, hash agreement, snapshot retention or finite check
establishes mathematical correctness.

No other model/worker/subagent, installation, external message,
source-cache edit, credential/auth/runtime/sandbox mutation,
campaign/private/SQL/transcript read or heavy compute was used. The
managed-runtime skill was read as required by higher-level environment
instructions; its status tool was unavailable and no credential
configuration was inspected. No supervised compute facility was exposed;
none was needed or launched.

Artifacts: P16.1-argument.md, P16.1-paper.tex, P16.1-paper.pdf,
n1-trace.md, n2-affine-tools.md, n3-runge-chain.md, n4-disposition.md,
canonical-notes.json, status.json,
source-extracts/source-provenance.json, external-premises.json,
finite-checks.json, command-log.json, build-pass-1.log and
build-pass-2.log. The unchanged mathematical target remains open in this
deliverable, candidate=false.

\section{Exact unchanged task and full
family}\label{exact-unchanged-task-and-full-family}

Original P16.1: P16.1 (T5): Finite-energy equations on singular or
unibranch varieties Extend the finite-energy non-Archimedean
Monge--Ampère operator and its existence-and- uniqueness theorem to
singular or unibranch varieties, replacing each use of smoothness in
compactness, envelopes, energy differentiation and comparison. State all
retained polarization and field hypotheses.

Complete unchanged mathematical family:

10 The general non-Archimedean Monge--Ampère equation Active items:
P16.1, P16.2, P16.3, P16.4, P16.5, P16.6, P16.7. Status: Open in the
residual scope.

10.1 Atomic targets P16.1 (T5): Finite-energy equations on singular or
unibranch varieties Extend the finite-energy non-Archimedean
Monge--Ampère operator and its existence-and- uniqueness theorem to
singular or unibranch varieties, replacing each use of smoothness in
compactness, envelopes, energy differentiation and comparison. State all
retained polarization and field hypotheses.

P16.2 (T5): Proper nonprojective algebraizable equation theories
Construct an appropriate replacement positivity datum and the associated
psh, finite-energy, comparison and equation theory on proper
nonprojective algebraizable spaces. A proper algebraic scheme with an
ample line bundle is already projective, so retaining that datum while
dropping projectivity is not a new case.

P16.3 (T5): Finite-energy equations in positive residue characteristic
Construct the required finite-energy operator and prove existence and
uniqueness when the original residue-characteristic restriction is
removed. Replace the resolution, model-regularity and envelope inputs
wherever they are used rather than merely changing the standing field
hypothesis.

P16.4 (T4): Finite-energy equations over arbitrary complete valued fields
Extend the operator and existence theorem to nondiscrete valuations and
arbitrary complete non-Archimedean fields outside the established
invariant-toric, skeleton-supported endpoint. A definition of the
finite-energy space alone does not establish the operator or the
equation.

P16.5 (T5): Finite-energy equations for nonample positivity classes
Extend the equation theory to big or pseudoeffective nonample classes,
specifying the psh cone, reference singularity type, energy and total
mass. At zero volume, the existing probability normalization must be
replaced rather than retained unchanged.

P16.6 (T4): Arbitrary-support finite-energy measures Prove the operator
and existence-and-uniqueness theorem for arbitrary finite-energy measure
support beyond the known toric skeletal regime. Measures must satisfy
the correct mass and energy conditions. Do not reduce arbitrary support
to a single skeleton or replace the genuine Monge--Ampère measure by
another inverse measure construction.

P16.7 (T5): Finite-energy equations on nonalgebraizable analytic spaces
Construct the finite-energy operator and its existence-and-uniqueness
theory on genuinely nonalgebraizable analytic spaces, with an explicit
positivity datum and the necessary com- pactness, envelope and
comparison replacements.

\begin{verbatim}
                                          73
\end{verbatim}

10.2 What constitutes a solution For each axis, state the retained
hypotheses and construct the finite-energy operator together with the
corresponding existence-and-uniqueness theorem. Specify the measure mass
and finite-energy condition, uniqueness modulo constants and the
replacement inputs for com- pactness, envelopes, energy differentiation
and comparison. An endpoint solving a special toric measure class does
not remove a general axis.

10.3 Full statement and detailed state of the art Problem 10.1. Extend
the finite-energy existence and uniqueness theory for the NA Monge--
Ampère equation beyond the currently proved combinations of smoothness,
projectivity, residue characteristic, valuation, positivity class,
support, and algebraizability hypotheses.

\begin{verbatim}
This is not one yes-or-no assertion. Positive characteristic, singular or unibranch varieties,
\end{verbatim}

nondiscrete valuations, big classes, arbitrary support, and
nonalgebraizable analytic spaces require different inputs. The known
variational proofs use compactness, envelope regularity,
differentiability of energy, and a comparison principle; deleting a
geometric hypothesis is valid only after replacing every place where it
enters. The relevant positive results and their distinct domains are
surveyed in Book Chapter 28, Book Chapter 29. The projectivity axis has
one logical constraint. A proper algebraic K-scheme endowed with an
ample line bundle is already projective over K {[}Stacksb, Tag 0B5N{]}.
Thus P16.2 cannot mean that one retains the BFJ line-bundle polarization
L while simply deleting the word ``projective.'' A nonprojective proper
algebraizable extension must first replace that positivity datum---and
then construct the corresponding psh, finite-energy, and comparison
theory---rather than merely rerun the original argument. No canonical
replacement with the required equation theory is asserted here. One
exact source-stated axis is especially clean. Pille--Schneider observes
that the finite- energy space can be defined over an arbitrary
non-Archimedean field, but that the extension of the Monge--Ampère
operator is not known there in general; the discretely valued
characteristic- zero and trivially valued cases are the established
endpoints {[}PS23, Remark 4.4{]}. This should be solved before bundling
all remaining axes into a single theorem.

Proposition 10.1 (Arbitrary-field invariant toric skeletal endpoint). Let
K be a complete nontrivially valued non-Archimedean field, XΣ a proper
toric K-variety, and L an ample toric line bundle. If µ is a positive
Radon measure on XΣan , is supported on the canonical skeleton NΣ ,
gives zero mass to NΣ \textbackslash NR , and has total mass degL (XΣ ),
then there is a toric continuous singular psh metric on L whose
Monge--Ampère measure is µ. It is unique up to multiplication of the
metric by a positive constant.

\begin{verbatim}
Thus nondiscreteness and the arbitrary complete ground field cause no obstruction in
\end{verbatim}

this invariant toric, skeleton-supported range. Here ``continuous
singular'' has the source's extended-valued meaning: local weights are
continuous into R ∪ \{−∞\}, as in {[}BGGJK25, §5.3.1{]}. It does not
imply finite real-valued weights or finite energy. This is not an
arbitrary- field finite-energy theorem in the BFJ sense: the result
permits a singular potential, requires torus invariance and skeletal
support, and does not construct the general extended Monge-- Ampère
operator. In particular, P16.4 and P16.6 retain the respectively
nonskeletal and noninvariant finite-energy residual scopes. There is also
a distinct, nonproper algebraic endpoint. Let n \textgreater{} d ≥ 1,
let X be a smooth projective Fano complex variety of dimension n, and
let D be a reduced simple-normal-crossing anticanonical divisor whose
dual complex is the standard (d−1)-simplex. For the affine Calabi-- Yau
complement X = X ~D, Wang's Theorem A constructs a continuous
non-Archimedean psh function on X an , over the trivially valued field C,
whose local Monge--Ampère measure

\begin{verbatim}
                                           74
\end{verbatim}

is the canonical Lebesgue measure on the essential skeleton {[}Wan25{]}.
The source is a version- one preprint and uses a local
CLD/tropical--convex construction. In particular, it neither constructs
the BFJ finite-energy operator nor treats arbitrary measures or arbitrary
affine varieties. Since X is affine and hence not proper, this endpoint does
not answer P16.2; it only records a special nonprojective algebraic
solution after the properness requirement is also dropped. There is a
separate finite-atomic inverse statement which must not be called a
finite-energy Monge--Ampère theorem. Under its standing assumptions that
K is maximally complete and \textbar K × \textbar{} = R\textgreater0 ,
Fang fixes an irreducible projective K-variety with an ample line bundle
and a finite set of graded birational points on its affine cone. For every
probability vector on that finite set, the source constructs a
Shilov-finite, possibly singular metric having the prescribed
equidistribution measure {[}Fan25, Proposition 4.8{]}. The source itself
says that comparing this Banach-algebra measure with the genuine
Monge--Ampère measure requires the separate con- tinuity property of an
envelope metric. Thus it supplies neither a continuous semipositive
solution nor the BFJ finite-energy operator for the prescribed atomic
measure, and does not narrow P16.6. There is also a resolved negative
regularity chronology. The first skeletal-factorization question of
Gubler--Jell--Künnemann--Martin already fails on their toric surface
{[}GJKM19, Appendix A{]}. Boucksom--Jonsson then construct a finite
atomic Monge--Ampère measure sup- ported on a dual complex whose
solution is invariant under no retraction attached to any SNC model;
they also construct a divisorial Green function that is not Q-PL
{[}BJ24b, Theo- rems A and B{]}. Hence support on a dual complex does
not imply factorization through some refined skeleton. This is a
theorem-level obstruction to a proof strategy, not an open regularity
expectation. The Rethlas programme treats each axis separately and first
searches the post-2020 liter- ature for an existing theorem. Alterations
and universal homeomorphisms can address some algebraic descent
questions, but they do not by themselves supply multiplier ideals or
dif- ferentiability in positive characteristic. The unrestricted
formulation is open; several strictly delimited axes are already known.

\begin{verbatim}
                                           75
\end{verbatim}

\subsection{Exact completion
criterion}\label{exact-completion-criterion}

Establish the genuine finite-energy operator, exact measure mass and
energy condition, existence, uniqueness modulo constants, and every
compactness/envelope/differentiation/comparison replacement for the
singular/geometrically unibranch projective domain with explicitly
retained field and polarization hypotheses. Verify the new envelope
source and the actual general-variety global theory. No change to the
original mathematical target is permitted.

\end{document}
